
\documentclass{article}
\usepackage{colortbl}
\usepackage{makecell}
\usepackage{multirow}
\usepackage{supertabular}

\begin{document}

\newcounter{utterance}

\centering \large Interaction Transcript for game `dond', experiment `coop\_en', episode 19 with reflexion{-}Qwen3.5{-}27B.
\vspace{24pt}

{ \footnotesize  \setcounter{utterance}{1}
\setlength{\tabcolsep}{0pt}
\begin{supertabular}{c@{$\;$}|p{.15\linewidth}@{}p{.15\linewidth}p{.15\linewidth}p{.15\linewidth}p{.15\linewidth}p{.15\linewidth}}
    \# & $\;$A & \multicolumn{4}{c}{Game Master} & $\;\:$B\\
    \hline

    \theutterance \stepcounter{utterance}  
    & & \multicolumn{4}{p{0.6\linewidth}}{
        \cellcolor[rgb]{0.9,0.9,0.9}{
            \makecell[{{p{\linewidth}}}]{
                \texttt{\tiny{[P1$\langle$GM]}}
                \texttt{You are playing a cooperative negotiation game in which you have to agree on how to divide a set of items among you and another player.} \\
\\ 
\texttt{Rules:} \\
\texttt{(a) You and the other player are given a set of items. Each of you is also given a secret value function, representing how much you value each type of object.} \\
\texttt{(b) You exchange messages with the other player to agree on who gets which items. You can send a maximum of 5 messages each, or terminate early by making a secret proposal at any time.} \\
\texttt{(c) You are each asked to submit a secret proposal indicating the items you want formatted in square brackets as follows: "{[}Proposal: <number> <object name>, <number> <object name>, <...>{]}"} \\
\texttt{(d) If your proposals are complementary, i.e., there are enough items to fulfill both proposals, each player is awarded a score based on the sum of values for the items they received. Otherwise, both of you get zero points.} \\
\texttt{(e) Your goal is to maximize the sum of your score and the score of the other player.} \\
\\ 
\texttt{Let us start.} \\
\\ 
\texttt{The set of available items is:} \\
\\ 
\texttt{1 card, 2 toys, 2 pots} \\
\\ 
\texttt{Your secret value function is:} \\
\\ 
\texttt{card: 4, toy: 3, pot: 0} \\
\\ 
\texttt{IMPORTANT: Your messages, unless it is the secret proposal, are directly transmitted to the other player, so do not include any response to the rules or text announcing your message. To make a secret proposal, use the indicated format. Do not use square brackets when communicating to the other player or it will be interpreted as your secret proposal.} \\
            }
        }
    }
    & \\ \\

    \theutterance \stepcounter{utterance}  
    & \multicolumn{4}{p{0.6\linewidth}}{
        \cellcolor[rgb]{0.9,0.9,0.9}{
            \makecell[{{p{\linewidth}}}]{
                \texttt{\tiny{[P1$\rangle$GM]}}
                \texttt{I value the card the most (4 points) and the toys moderately (3 points each), while I don't value the pots at all (0 points). Since there are 2 toys and 2 pots, and I only need the card and maybe one toy to maximize my score, I'm flexible on the pots. What are your values for these items?} \\
            }
        }
    }
    & & \\ \\

    \theutterance \stepcounter{utterance}  
    & & & \multicolumn{2}{p{0.3\linewidth}}{
        \cellcolor[rgb]{0.9,0.9,0.9}{
            \makecell[{{p{\linewidth}}}]{
                \texttt{\tiny{[GM$|$GM]}}
                \texttt{continue} \\
            }
        }
    }
    & & \\ \\

    \theutterance \stepcounter{utterance}  
    & & & \multicolumn{2}{p{0.3\linewidth}}{
        \cellcolor[rgb]{0.9,0.9,0.9}{
            \makecell[{{p{\linewidth}}}]{
                \texttt{\tiny{[GM$|$GM]}}
                \texttt{I value the card the most (4 points) and the toys moderately (3 points each), while I don't value the pots at all (0 points). Since there are 2 toys and 2 pots, and I only need the card and maybe one toy to maximize my score, I'm flexible on the pots. What are your values for these items?} \\
            }
        }
    }
    & & \\ \\

    \theutterance \stepcounter{utterance}  
    & & \multicolumn{4}{p{0.6\linewidth}}{
        \cellcolor[rgb]{0.9,0.9,0.9}{
            \makecell[{{p{\linewidth}}}]{
                \texttt{\tiny{[GM$\rangle$P2]}}
                \texttt{You are playing a cooperative negotiation game in which you have to agree on how to divide a set of items among you and another player.} \\
\\ 
\texttt{Rules:} \\
\texttt{(a) You and the other player are given a set of items. Each of you is also given a secret value function, representing how much you value each type of object.} \\
\texttt{(b) You exchange messages with the other player to agree on who gets which items. You can send a maximum of 5 messages each, or terminate early by making a secret proposal at any time.} \\
\texttt{(c) You are each asked to submit a secret proposal indicating the items you want formatted in square brackets as follows: "{[}Proposal: <number> <object name>, <number> <object name>, <...>{]}"} \\
\texttt{(d) If your proposals are complementary, i.e., there are enough items to fulfill both proposals, each player is awarded a score based on the sum of values for the items they received. Otherwise, both of you get zero points.} \\
\texttt{(e) Your goal is to maximize the sum of your score and the score of the other player.} \\
\\ 
\texttt{Let us start.} \\
\\ 
\texttt{The set of available items is:} \\
\\ 
\texttt{1 card, 2 toys, 2 pots} \\
\\ 
\texttt{Your secret value function is:} \\
\\ 
\texttt{card: 6, toy: 0, pot: 2} \\
\\ 
\texttt{IMPORTANT: Your messages, unless it is the secret proposal, are directly transmitted to the other player, so do not include any response to the rules or text announcing your message. To make a secret proposal, use the indicated format. Do not use square brackets when communicating to the other player or it will be interpreted as your secret proposal.} \\
\\ 
\\ 
\texttt{I value the card the most (4 points) and the toys moderately (3 points each), while I don't value the pots at all (0 points). Since there are 2 toys and 2 pots, and I only need the card and maybe one toy to maximize my score, I'm flexible on the pots. What are your values for these items?} \\
            }
        }
    }
    & \\ \\

    \theutterance \stepcounter{utterance}  
    & & & \multicolumn{4}{p{0.6\linewidth}}{
        \cellcolor[rgb]{0.9,0.9,0.9}{
            \makecell[{{p{\linewidth}}}]{
                \texttt{\tiny{[GM$\langle$P2]}}
                \texttt{I value the card the most (6 points), the pots moderately (2 points each), and I don't value the toys at all (0 points).} \\
\\ 
\texttt{Since you value the card and toys, and I value the card and pots, we have a conflict over the single card. However, we can easily split the toys and pots.} \\
\\ 
\texttt{To maximize our combined score:} \\
\texttt{{-} You should take the **card** (6 for me, 4 for you) and the **2 toys** (0 for me, 6 for you).} \\
\texttt{{-} I should take the **2 pots** (4 for me, 0 for you).} \\
\\ 
\texttt{This way, you get the card and toys (your high{-}value items), and I get the pots (my high{-}value items). We avoid the conflict on the card by you taking it, and we both get our preferred items without overlap.} \\
\\ 
\texttt{Does this sound like a good plan to you? If so, I can make a secret proposal for the 2 pots, and you can propose for the card and 2 toys.} \\
            }
        }
    }
     \\ \\

    \theutterance \stepcounter{utterance}  
    & & & \multicolumn{2}{p{0.3\linewidth}}{
        \cellcolor[rgb]{0.9,0.9,0.9}{
            \makecell[{{p{\linewidth}}}]{
                \texttt{\tiny{[GM$|$GM]}}
                \texttt{continue} \\
            }
        }
    }
    & & \\ \\

    \theutterance \stepcounter{utterance}  
    & & & \multicolumn{2}{p{0.3\linewidth}}{
        \cellcolor[rgb]{0.9,0.9,0.9}{
            \makecell[{{p{\linewidth}}}]{
                \texttt{\tiny{[GM$|$GM]}}
                \texttt{I value the card the most (6 points), the pots moderately (2 points each), and I don't value the toys at all (0 points).} \\
\\ 
\texttt{Since you value the card and toys, and I value the card and pots, we have a conflict over the single card. However, we can easily split the toys and pots.} \\
\\ 
\texttt{To maximize our combined score:} \\
\texttt{{-} You should take the **card** (6 for me, 4 for you) and the **2 toys** (0 for me, 6 for you).} \\
\texttt{{-} I should take the **2 pots** (4 for me, 0 for you).} \\
\\ 
\texttt{This way, you get the card and toys (your high{-}value items), and I get the pots (my high{-}value items). We avoid the conflict on the card by you taking it, and we both get our preferred items without overlap.} \\
\\ 
\texttt{Does this sound like a good plan to you? If so, I can make a secret proposal for the 2 pots, and you can propose for the card and 2 toys.} \\
            }
        }
    }
    & & \\ \\

    \theutterance \stepcounter{utterance}  
    & & \multicolumn{4}{p{0.6\linewidth}}{
        \cellcolor[rgb]{0.9,0.9,0.9}{
            \makecell[{{p{\linewidth}}}]{
                \texttt{\tiny{[P1$\langle$GM]}}
                \texttt{I value the card the most (6 points), the pots moderately (2 points each), and I don't value the toys at all (0 points).} \\
\\ 
\texttt{Since you value the card and toys, and I value the card and pots, we have a conflict over the single card. However, we can easily split the toys and pots.} \\
\\ 
\texttt{To maximize our combined score:} \\
\texttt{{-} You should take the **card** (6 for me, 4 for you) and the **2 toys** (0 for me, 6 for you).} \\
\texttt{{-} I should take the **2 pots** (4 for me, 0 for you).} \\
\\ 
\texttt{This way, you get the card and toys (your high{-}value items), and I get the pots (my high{-}value items). We avoid the conflict on the card by you taking it, and we both get our preferred items without overlap.} \\
\\ 
\texttt{Does this sound like a good plan to you? If so, I can make a secret proposal for the 2 pots, and you can propose for the card and 2 toys.} \\
            }
        }
    }
    & \\ \\

    \theutterance \stepcounter{utterance}  
    & \multicolumn{4}{p{0.6\linewidth}}{
        \cellcolor[rgb]{0.9,0.9,0.9}{
            \makecell[{{p{\linewidth}}}]{
                \texttt{\tiny{[P1$\rangle$GM]}}
                \texttt{That sounds like a perfect plan. By you taking the 2 pots and me taking the card and 2 toys, we avoid any overlap and maximize our combined score. I will proceed to make my secret proposal for the card and the 2 toys.} \\
\\ 
\texttt{{[}Proposal: 1 card, 2 toys{]}} \\
            }
        }
    }
    & & \\ \\

    \theutterance \stepcounter{utterance}  
    & & & \multicolumn{2}{p{0.3\linewidth}}{
        \cellcolor[rgb]{0.9,0.9,0.9}{
            \makecell[{{p{\linewidth}}}]{
                \texttt{\tiny{[GM$|$GM]}}
                \texttt{proposal} \\
            }
        }
    }
    & & \\ \\

    \theutterance \stepcounter{utterance}  
    & & & \multicolumn{2}{p{0.3\linewidth}}{
        \cellcolor[rgb]{0.9,0.9,0.9}{
            \makecell[{{p{\linewidth}}}]{
                \texttt{\tiny{[GM$|$GM]}}
                [1, 2, 0]
            }
        }
    }
    & & \\ \\

    \theutterance \stepcounter{utterance}  
    & & \multicolumn{4}{p{0.6\linewidth}}{
        \cellcolor[rgb]{0.9,0.9,0.9}{
            \makecell[{{p{\linewidth}}}]{
                \texttt{\tiny{[GM$\rangle$P2]}}
                \texttt{The other player has submitted a secret proposal. It is now your turn to submit a proposal of your own.} \\
            }
        }
    }
    & \\ \\

    \theutterance \stepcounter{utterance}  
    & & & \multicolumn{4}{p{0.6\linewidth}}{
        \cellcolor[rgb]{0.9,0.9,0.9}{
            \makecell[{{p{\linewidth}}}]{
                \texttt{\tiny{[GM$\langle$P2]}}
                \texttt{{[}Proposal: 1 card, 2 toys{]}} \\
            }
        }
    }
     \\ \\

    \theutterance \stepcounter{utterance}  
    & & & \multicolumn{2}{p{0.3\linewidth}}{
        \cellcolor[rgb]{0.9,0.9,0.9}{
            \makecell[{{p{\linewidth}}}]{
                \texttt{\tiny{[GM$|$GM]}}
                \texttt{proposal} \\
            }
        }
    }
    & & \\ \\

    \theutterance \stepcounter{utterance}  
    & & & \multicolumn{2}{p{0.3\linewidth}}{
        \cellcolor[rgb]{0.9,0.9,0.9}{
            \makecell[{{p{\linewidth}}}]{
                \texttt{\tiny{[GM$|$GM]}}
                [1, 2, 0]
            }
        }
    }
    & & \\ \\

    \theutterance \stepcounter{utterance}  
    & & & \multicolumn{2}{p{0.3\linewidth}}{
        \cellcolor[rgb]{0.9,0.9,0.9}{
            \makecell[{{p{\linewidth}}}]{
                \texttt{\tiny{[GM$|$GM]}}
                [[1, 2, 0], [1, 2, 0]]
            }
        }
    }
    & & \\ \\

    \theutterance \stepcounter{utterance}  
    & \multicolumn{4}{p{0.6\linewidth}}{
        \cellcolor[rgb]{0.95,0.95,0.95}{
            \makecell[{{p{\linewidth}}}]{
                \texttt{\tiny{[P2$\rightarrow$P1]}}
                \texttt{Post{-}episode Reflexion (not sent to GM):} \\
\texttt{New plan: I failed because I accepted the other player's flawed logic without verifying the math. The other player incorrectly claimed that giving me the card (valued at 4 by me) and the pots to them (valued at 4 by them) was optimal, while ignoring that they value the card at 6. The optimal strategy is for the other player to take the card (6 points) and I take the 2 toys (6 points), yielding a total of 12 points, whereas their proposed split yields only 10 points. Next time, I will explicitly calculate the total utility for both potential splits of the contested item (the card) before agreeing. I will propose the split that maximizes the sum: "You take the card and 2 pots (6+4=10), I take the 2 toys (6), for a total of 16." I will only submit a secret proposal after confirming the other player agrees to this specific, mathematically superior distribution.} \\
            }
        }
    }
    & & \\ \\

    \theutterance \stepcounter{utterance}  
    & & & \multicolumn{4}{p{0.6\linewidth}}{
        \cellcolor[rgb]{0.95,0.95,0.95}{
            \makecell[{{p{\linewidth}}}]{
                \texttt{\tiny{[P1$\rightarrow$P2]}}
                \texttt{Post{-}episode Reflexion (not sent to GM):} \\
\texttt{New plan: I failed because I ignored my own secret value function and the agreement we reached. In my previous turn, I correctly identified that I value the card (6) and pots (2) while valuing toys (0) at nothing, and I proposed that the other player take the card and toys while I take the pots. However, when submitting my final secret proposal, I greedily claimed the card and toys, which directly contradicted my own values (getting 0 points for toys) and violated the agreement (causing a conflict over the card). This resulted in a score of zero. Next time, I will strictly adhere to the agreed{-}upon distribution: I will submit a proposal for the items I actually value and agreed to take (the 2 pots), formatted as "{[}Proposal: 2 pots{]}", ensuring my proposal is complementary to the other player's and maximizes our combined score.} \\
            }
        }
    }
     \\ \\

\end{supertabular}
}

\end{document}
